\documentclass[11pt]{article}
\usepackage[margin=1.15in]{geometry}
\usepackage{amsmath,amssymb,graphicx,xcolor,booktabs,needspace}
\usepackage[hidelinks]{hyperref}
\newcommand{\Cov}{\mathsf{C}}
\newcommand{\M}{\mathsf{M}}
\newcommand{\I}{\mathsf{I}}
\newcommand{\T}{^{\mathsf{T}}}
\newcommand{\E}{\mathbb{E}}
\newcommand{\Var}{\operatorname{Var}}
\newcommand{\cov}{\operatorname{Cov}}
\emergencystretch=3em
\setlength{\parskip}{0.2em}
\begin{document}
\begin{center}
{\LARGE Learning a covariance from simulations\par}
\vspace{0.5em}
{\large The Friedrich--Eifler precision expansion and Rafael's DES shear maps\par}
\vspace{0.5em}
20 November 2026
\end{center}

\section{The question and how to read this note}
\paragraph{Why simulate a survey more than once?}
Imagine measuring the same cosmological model in many different Universes.
The galaxies and dark-matter structures would be in different places, so our
measurement would change even if the underlying parameters did not.
Galaxy shapes and measurement noise introduce further changes. A covariance
matrix describes the size and correlation of those fluctuations. We need it
because a difference between theory and data is meaningful only in relation
to how much the data could fluctuate.

\paragraph{What we want to improve.}
An analytic covariance is quick to evaluate and has no Monte Carlo sampling
noise. Its physical approximations, however, can miss correlations produced
by nonlinear structure, the survey boundary or the measurement procedure.
Simulated surveys contain some of these effects, but estimating a large
matrix from a finite set is noisy. Friedrich and Eifler~\cite{pme} proposed
expanding the \emph{inverse covariance} around an analytic model and learning
the correction from simulations. This note derives that idea before asking
how Rafael Gomes's downloaded shear maps could supply the measurements.
Improvement means more reliable uncertainties; it need not mean smaller ones.

\paragraph{A route through the mathematics.}
Sections 2--4 introduce the measurement, the maps and the noisy-matrix
problem. Sections 5--6 derive the expansion, beginning with one number rather
than a matrix. Section 7 explains a complication specific to repeated skies
and noise draws. Sections 8--9 give a limited pilot and its validation criteria.
An understanding of averages, variances and matrix multiplication is enough
to follow the argument; the Gaussian matrix-moment identity is stated with
its assumptions. The companion notebook first uses artificial Gaussian
vectors with known covariance, then describes the real-map workflow. Its
existing numerical results are not a calibration of the downloaded maps.

\paragraph{Credits and present scope.}
\textbf{We thank Rafael Gomes for supplying the maps and explaining their
pixel selection, and Sunao Sugiyama and Rafael Gomes for their theory and
measurement software.} The relevant analysis is Gomes et al.~\cite{gomes};
the statistical method is due to Oliver Friedrich and Tim Eifler.
The local download is incomplete. We inspected two array headers and small
slices, not the complete collection. We have not measured its joint vectors,
corrected a DES covariance or run a new Fisher forecast. This is an explanation
and a plan that can be checked before undertaking those computations.

\section{What the covariance is a covariance of}
\paragraph{The observable comes before its uncertainty.}
Weak gravitational lensing distorts the apparent shapes of distant galaxies.
Shear has two components because an elongated image has both an amplitude
and an orientation. A shear map stores these components at sky positions.
Our inference does not use each pixel as a separate observable. Instead, it
uses averages of products of shears: two-point correlations compare pairs
of positions, while third-order statistics compare triples. Aperture mass,
written $M_{\rm ap}$, is a filtered lensing statistic; its third moment is an
average of three such filtered measurements. We abbreviate that observable
as Map3, not as a third independent shear component.

\paragraph{A vector makes the bookkeeping explicit.}
Let $\boldsymbol d$ be a column of $p$ measured numbers. Its entries include
$\xi_+$ and $\xi_-$, the two usual shear correlation functions, and Map3,
with their specified angular bins and source-redshift combinations. In
Rafael's methods setup, 400 two-point entries and 80 aperture entries give
480 entries before further selections~\cite{gomes}. This count describes that
setup, not an obligation to use every row here. Changing a scale cut changes
the vector and therefore changes the covariance we must estimate.
Every row needs a definition, including its units, bins, weights and cuts.

\paragraph{A covariance records shared fluctuations.}
Write $\boldsymbol\mu=\E[\boldsymbol d]$ for the ensemble mean. Then
\begin{equation}
 \Cov_{ij}=\E[(d_i-\mu_i)(d_j-\mu_j)].
\end{equation}
The diagonal gives individual variances. A positive off-diagonal entry means
that the two observables tend to rise together. For example, a sky containing
an unusually large concentration of matter can change both two- and
three-point statistics. In block form,
\begin{equation}
 \Cov=\begin{pmatrix}\Cov_{22}&\Cov_{23}\\
                     \Cov_{23}\T&\Cov_{33}\end{pmatrix}.
\end{equation}
The block $\Cov_{23}$ is why the two statistics must be measured on the
\emph{same} simulated survey before their vectors are concatenated. Drawing
unrelated skies for the two blocks erases the shared fluctuations.

\paragraph{Why the inverse matters.}
The precision matrix is $\Psi=\Cov^{-1}$. For a Gaussian likelihood with a
fixed covariance, the squared distance of a model $\boldsymbol m$ from the
data is
\begin{equation}
 \chi^2=(\boldsymbol d-\boldsymbol m)\T\Psi
                  (\boldsymbol d-\boldsymbol m).
\end{equation}
Consider two observables with covariance
$\sigma^2\left(\begin{smallmatrix}1&r\\r&1\end{smallmatrix}\right)$,
where $\sigma$ is their common standard deviation and $r$ their correlation.
The normalized sum and difference have variances $\sigma^2(1+r)$ and
$\sigma^2(1-r)$. If $r$ is close to one, their difference is tightly measured.
An error in $r$ then matters greatly for its inverse variance. This is why
judging only the diagonal entries of a covariance is insufficient.

\section{What is in the downloaded folder}
\paragraph{Confirmed information.}
Rafael's message identifies the files as CosmoGrid shear maps at HEALPix
resolution $N_{\rm side}=512$. They contain only pixels where the DES
footprint is nonzero; the accompanying \texttt{indices.npy} lists these sky
pixels. Two files, with labels \texttt{run\_0000}, \texttt{ncut0} and
\texttt{nreal0} or \texttt{nreal1}, have shape $(4,2,369190)$ and double
precision entries. Each file contains 23,628,288 bytes including its header.
The first 4,096 pixels in each slot were finite. This small check establishes
neither completeness of the download nor its scientific calibration.

\paragraph{Array positions are not sky coordinates.}
The likely first two axes are four source-redshift bins and two shear
components; their ordering and signs still require confirmation. If the
index file is a common one-dimensional list, the entry $d_{bgj}$ in a map
belongs to HEALPix pixel \texttt{indices[j]}. The index $j$ is merely its
position in the compressed list. Sorting that list without sorting the
shear values in exactly the same way places the signal at incorrect sky
positions. A rectangular reshape is also incorrect. Pixels outside the list
are unobserved, not measurements of zero shear.

\paragraph{Recovering a footprint picture.}
HEALPix divides the sphere into $12N_{\rm side}^2$ equal-area pixels. For 512
this is 3,145,728 pixels. There are two index orderings, RING and NESTED,
with the same valid integer range but different coordinate assignments.
The producer must identify which one applies, and whether its coordinates
are equatorial. The viewing notebook uses \texttt{healpy.pix2ang} after those
choices are supplied. At writing, \texttt{indices.npy} was not present locally,
so we have not shown a guessed footprint. The notebook also checks whether
the indices are a common list; a separate list for each cut needs an explicit
selection of the matching list.

\paragraph{A preview is not a measurement.}
The viewing notebook samples at most 30,000 retained positions from each of
two named files and plots both components of the first source-bin slot.
This keeps rendering inexpensive. The same color limits permit visual
comparison, but no quantitative estimator should use this display
subsampling. A component image also does not validate the spin-2 coordinate
basis: a shear correlation estimator rotates components relative to pair
or triangle directions and needs the correct sign convention. Pixel weights,
galaxy occupancies and the source distributions are further measurement
inputs; pixel indices alone do not supply them.

\section{Why simply inverting the sample covariance is difficult}
\paragraph{The familiar estimator.}
For $N$ independent measurements of the same vector, the unbiased sample
covariance is
\begin{equation}
 \widehat\Cov=\frac{1}{N-1}\sum_{r=1}^{N}
 (\boldsymbol d_r-\bar{\boldsymbol d})
 (\boldsymbol d_r-\bar{\boldsymbol d})\T,
 \qquad \bar{\boldsymbol d}=\frac1N\sum_r\boldsymbol d_r.
\end{equation}
The $N-1$ accounts for estimating the mean from these measurements. After
subtracting it, the residual vectors sum to zero, so there can be at most
$N-1$ independent directions. Consequently $\widehat\Cov$ has rank at most
$N-1$. A 480-component vector and 199 independent measurements give rank at
most 198: its full sample covariance cannot be inverted.

\paragraph{Even an invertible estimate fluctuates.}
For a single positive variance estimate $\hat v=v+\epsilon$, expand its
inverse for small $\epsilon/v$:
\begin{equation}
 \E[\hat v^{-1}]\simeq v^{-1}
 +\frac{\E[\epsilon^2]}{v^3},\qquad \E[\epsilon]=0.
\end{equation}
Zero average error in the variance therefore does not give zero average
error in its inverse. Matrix inversion has the same problem and mixes errors
between entries. A correction to the mean bias of an inverse, such as the
Hartlap factor under its Gaussian assumptions, does not eliminate its
sampling fluctuations or make a singular matrix invertible.

\paragraph{The analytic model supplies information.}
Let $\M$ be a positive-definite analytic approximation to $\Cov$. PME only
inverts $\M$, so a rank-deficient sample covariance is allowed. This does not
create missing simulation information. Directions poorly constrained by the
simulations depend strongly on the analytic baseline. We must validate its
physical ingredients and compare corrections in the parameter combinations
we actually use. Nor can one transfer a simulation-count requirement from
another forecast without checking its vector, baseline and precision target.

\section{Deriving the precision expansion}
\paragraph{Start with one dimension.}
If the modeled variance is one and the true variance is $1+x$, then
\begin{equation}
 \frac1{1+x}=1-x+x^2-x^3+\cdots,\qquad |x|<1.
\end{equation}
At $x=0.1$, the true inverse is $0.90909\ldots$; the first- and second-order
approximations are $0.9$ and $0.91$. The expansion works because the
\emph{fractional error} is small. The same absolute variance error could be
large or small depending on the variance used to scale it.

\paragraph{Put different observables on comparable scales.}
Two- and three-point entries can have very different numerical sizes.
Factor the model as $\M=L L\T$, where $L$ is its lower-triangular Cholesky
factor. Transform a residual to $\boldsymbol y=L^{-1}(\boldsymbol d-\boldsymbol\mu)$.
Its modeled covariance is the identity. This operation, called whitening,
removes the model's units and correlations. The true covariance in these
coordinates is
\begin{equation}
 R=L^{-1}\Cov L^{-\mathsf T}=\I+X,
 \qquad X=L^{-1}(\Cov-\M)L^{-\mathsf T}.
\end{equation}
Here $L^{-\mathsf T}$ means $(L\T)^{-1}$ and $X$ is the dimensionless model
error. Whitening can be evaluated by triangular solves rather than by
forming explicit numerical inverses.

\paragraph{Apply the scalar argument to each matrix direction.}
A real symmetric $X$ has orthogonal eigenvectors. Along an eigenvector with
eigenvalue $x$, the required inverse factor is precisely $1/(1+x)$.
Therefore
\begin{equation}
 \Psi=L^{-\mathsf T}(\I-X+X^2-\cdots)L^{-1}.
 \label{eq:series}
\end{equation}
If $\rho=\max|x|<1$, the norm of the remainder after second order is at most
$\rho^3/(1-\rho)$ in whitened coordinates. This is a statement about the
\emph{true} model error. A sample estimate of $X$ is noisy. In particular,
a singular whitened sample covariance has zero eigenvalues, so subtracting
$\I$ produces eigenvalues of $-1$ even when the true model is excellent.
The sample spectral radius is useful to record, but is neither a proof nor
a general disproof of true convergence with few simulations.

\paragraph{Truncation and simulation noise are different errors.}
Even an infinite simulation ensemble leaves the discarded cubic and higher
terms when using a second-order formula. A finite ensemble adds random
errors in the retained terms. Improving the model can reduce truncation;
adding independent skies can reduce sampling errors. These operations are
not interchangeable. A tiny pilot tests file interpretation and algebra,
not convergence of a high-dimensional precision estimate.

\section{Why the noisy square must be corrected}
\paragraph{A square does not average like its argument.}
Let $S=L^{-1}\widehat\Cov L^{-\mathsf T}$ be the whitened sample covariance.
Its expectation is $R$, but generally $\E[S^2]\ne R^2$. Write $S=R+\delta S$;
then the average square contains $\E[\delta S\delta S]$ as well as $R^2$.
Simply inserting the noisy residual $S-\I$ into Equation~\eqref{eq:series}
therefore introduces a second-order sampling bias.

\paragraph{The Gaussian identity and what it assumes.}
For independent Gaussian vectors and their ordinary mean-subtracted sample
covariance, $\nu S$ follows a Wishart distribution with $\nu=N-1$ degrees of
freedom and scale $R$. Wishart is the matrix analogue of the distribution
of a sum of squared Gaussian numbers. Its moment identities give
\begin{align}
 \E[S^2]&=(1+1/\nu)R^2+R\,\operatorname{tr}(R)/\nu,\\
 \E[S\operatorname{tr}(S)]&=R\operatorname{tr}(R)+2R^2/\nu.
\end{align}
The trace, $\operatorname{tr}(R)$, is the sum of diagonal entries. Combining
these relations cancels the extra trace term:
\begin{equation}
 Q=\frac{\nu^2S^2-\nu S\operatorname{tr}(S)}{\nu^2+\nu-2},
 \qquad \E[Q]=R^2,\qquad N>2.
\end{equation}
Thus the first- and second-order precision estimators are
\begin{align}
 \widehat\Psi_1&=L^{-\mathsf T}(2\I-S)L^{-1},\\
 \widehat\Psi_2&=L^{-\mathsf T}(3\I-3S+Q)L^{-1}.
\end{align}
Their expectations equal the corresponding truncated expansions, not the
exact precision. They need no Hartlap multiplier. A realization can be
indefinite; negative precision eigenvalues must be diagnosed rather than
clipped away and silently used in a likelihood.

\paragraph{What if the aperture vectors are not Gaussian?}
Nonlinear summaries need not follow a multivariate Gaussian distribution.
An available alternative for the quadratic expectation is to estimate
covariances $S_1,S_2$ from two \emph{independent} parent groups measuring the
same population. Independence gives
\begin{equation}
 \E[(S_1S_2+S_2S_1)/2]=R^2.
\end{equation}
Replacing $S$ by $(S_1+S_2)/2$ and $Q$ by this symmetric product in the
second-order formula removes the Gaussian assumption for that expectation.
It still requires finite moments, unbiased sample covariances and a model
fixed independently of both groups. It does not remove truncation or
sampling noise; splitting a small ensemble can make the latter large.
A third independent group is needed for a clean held-out check.

\paragraph{The component version can be better, when its inputs exist.}
The especially useful case in Friedrich--Eifler is
$\Cov=A+B$ with a well-known covariance contribution $A$ and model
$\M=A+B_m$. If simulations directly isolate $B$, they need not re-estimate
$A$ with Monte Carlo noise. Define $S_B=L^{-1}\widehat B L^{-\mathsf T}$,
$T=L^{-1}B_m L^{-\mathsf T}$, and let $Q_B$ be the above unbiased-square
expression applied to $S_B$. The second-order polynomial is
\begin{equation}
 \I-S_B+T+Q_B-S_BT-TS_B+T^2.
\end{equation}
Both multiplication orders matter. These are additive covariance
components, \emph{not} the two-point and three-point blocks with their
cross-covariance discarded. Nor does subtracting a known $A$ from a full
sample covariance turn the remainder into a Wishart sample of $B$.
The existing component routine must receive actual isolated-component
realizations satisfying its assumptions.

\section{Repeated noise is useful, but it is not a new sky}
\paragraph{Three labels need three meanings.}
The filenames contain \texttt{run}, \texttt{ncut} and \texttt{nreal}.
They suggest parent sky, footprint cut and noise realization; Rafael's
message confirms the pixel selection, not this interpretation. His methods
paper describes 199 parent skies and four DES footprints per sky, giving
796 mocks~\cite{gomes}. That description alone does not identify every
relationship in this new archive. Obtain its parent IDs and noise-generation
conventions before assigning a statistical degree of freedom. Keep all
children of one parent together when splitting training and validation.

\paragraph{A useful decomposition from conditional averages.}
Suppose, for this derivation, a fixed sky $s$ has repeated independent noise
realizations. Let $\boldsymbol d=\mathcal T(s+n)$ be the measured vector,
where $n$ denotes shape noise and $\mathcal T$ the correlation measurement.
Define $\boldsymbol\mu_s=\E_n[\boldsymbol d\mid s]$. Add and subtract this
conditional mean in $\boldsymbol d-\E[\boldsymbol d]$. The cross term has
zero average because $\E_n[\boldsymbol d-\boldsymbol\mu_s\mid s]=0$. Hence
\begin{equation}
 \Cov=\underbrace{\cov_s(\boldsymbol\mu_s)}_{B}
       +\underbrace{\E_s[\cov_n(\boldsymbol d\mid s)]}_{A}.
 \label{eq:total}
\end{equation}
The first term describes changing skies; the second describes changing
noise on a fixed sky, averaged over skies. This identity requires no Gaussian
assumption. Its physical usefulness depends on what random quantities the
simulation actually varies or holds fixed.

\paragraph{Averaging reduces noise, which changes the target.}
Take $K>1$ noise draws for each of $R_s>1$ independent parents, using one
fixed footprint selection and the same measurement throughout. Let
$\bar{\boldsymbol d}_s$ be the mean vector for parent $s$, and $W_s$ the
sample covariance of its $K$ vectors with denominator $K-1$. Define
$\widehat A=R_s^{-1}\sum_s W_s$ and let $\widehat C_{\bar d}$ be the ordinary
sample covariance of the $R_s$ parent means. Conditional independence and
equal $K$ imply
\begin{equation}
 \E[\widehat A]=A,\qquad
 \E[\widehat C_{\bar d}]=B+A/K.
\end{equation}
An unbiased estimate of the covariance of one noisy survey is consequently
\begin{equation}
 \widehat\Cov=\widehat C_{\bar d}+(1-1/K)\widehat A,
 \qquad
 \widehat B=\widehat C_{\bar d}-\widehat A/K.
 \label{eq:hierarchical}
\end{equation}
For example, averaging 20 noise realizations leaves only one twentieth of
the conditional-noise covariance in the parent means. Using their covariance
alone would miss the other nineteen twentieths of the single-observation
noise. This is a worked statistical example, not a claim that the local
\texttt{nreal} labels have already been verified as these 20 draws.

\paragraph{This decomposition does not yet supply component PME.}
The difference defining $\widehat B$ in Equation~\eqref{eq:hierarchical}
can have negative eigenvalues and is not a Wishart sample covariance.
Inserting it into the component formula with an invented effective sample
count is unjustified. Also, $\mathcal T$ is nonlinear: measuring an averaged
map is generally different from averaging measured vectors. Mixed signal
and noise terms enter the covariance of products, so noiseless maps do not
automatically deliver the $B$ needed for a particular proposed split.
Repeated-noise data can help test those terms, but deriving the required
product-moment correction is a separate statistical task.

\paragraph{Footprint cuts require another check.}
Different cuts from the same full sky can share large-scale modes or even
pixels. Averaging them changes the effective survey rather than producing
one ordinary DES realization. A conservative initial construction is one
preselected cut and one preselected noise draw per independent parent.
This uses less of the archive but gives a clear reference population.
Using more cuts requires measuring their correlations or deriving an
estimator that accounts for them. Parent-level resampling can be useful
for uncertainty checks, but does not magically make cuts independent.

\section{The measurement must match before the covariance can match}
\paragraph{Pixel size is a physical smoothing scale.}
A HEALPix pixel has area $4\pi/(12N_{\rm side}^2)$ steradians. Its square-root
angular scale is about $6.87$ arcminutes at 512, twice the scale at 1024.
The methods paper uses 1024 and discusses a small-triangle cutoff near
8 arcminutes~\cite{gomes}. At 512 the same numerical cut is much closer to
a pixel width. The old cutoff is therefore not a validation of the new maps.
Increasing it in a sensitivity test is sensible, but no factor-of-two rule
by itself establishes a converged estimator. Small aperture radii also need
an explicit resolution check.

\paragraph{A mask is more than an area.}
Two footprints with the same total area can have different pair and triangle
counts, holes and boundaries. A pixel-weighted map estimator may also differ
from a galaxy-catalog estimator. Record source-bin definitions, pixel weights,
shape-noise prescription, coordinate basis, pair/triangle bins and angular
averaging, aperture filters, finite-triangle cuts and any pixel smoothing.
The analytic covariance and the simulated measurement must target these
same choices. A mismatch in their means can reveal a problem, but agreement
of means alone does not establish agreement of covariances.

\paragraph{Physics absent from the maps cannot be calibrated by them.}
The current message does not specify intrinsic alignment, baryonic feedback,
source clustering, calibration fluctuations or the cosmological parameters
of the archive. A gravitational shear ensemble at a fixed cosmology cannot
by itself validate all these contributions or the covariance's cosmology
dependence. Confirm which ingredients are simulated, which are modeled
analytically and which are neglected. The current full-filter gravitational
covariance implementation is not already a matched finite-triangle TATT
covariance for this download. Correcting a matrix cannot repair a wrong
mean model or supply a physical effect absent from both model and simulations.

\section{A bounded pilot and the decision it should support}
\paragraph{First establish the inputs.}
Obtain \texttt{indices.npy}, the ordering, coordinate frame and label meanings;
verify two complete files against the producer's expected size or checksum.
Make the footprint previews, then measure one joint vector with fully
recorded rows and settings. Repeat with one more selected map to check
pairing and repeatability. The first purpose is to catch a wrong coordinate
assignment, sign or row order cheaply. Do not begin by passing the entire
partially downloaded directory to a measurement job.

\paragraph{Then construct a defensible training set.}
Choose parents before inspecting their covariance residuals. Reserve separate
parents for model tuning, correction estimation and validation. Use one
specified footprint/noise draw per parent for the initial full-covariance
route. Store generated vectors and their labels under \texttt{chains/}.
The existing notebook loader checks row ordering and rejects repeated or
excluded parent IDs, but the analyst must still establish their physical
independence. If two independent groups are used for the quadratic product,
keep validation parents out of both. A small pilot may simply be too noisy
to support any correction; that is an informative result.

\paragraph{Compare at fixed scientific assumptions.}
Evaluate the analytic precision, first-order PME and, if justified,
second-order or split-group PME. Keep the mean, data rows, derivatives and
priors fixed. Record precision eigenvalues, corrections to correlations,
and parameter uncertainties. For parameter derivatives $D_{i\alpha}
=\partial m_i/\partial\theta_\alpha$, the Gaussian fixed-covariance Fisher
matrix is $F=D\T\Psi D+F_{\rm prior}$. Here $F_{\rm prior}$ represents the
specified prior information. This projects the covariance changes onto
parameter-sensitive directions. A higher figure of merit merely means a
smaller quoted ellipse; it does not establish more accurate error bars.

\paragraph{Validate weights using data not used to learn them.}
For independent residuals $\boldsymbol r=\boldsymbol d-\boldsymbol\mu$
about a known ensemble mean,
\begin{equation}
 \E[\boldsymbol r\T\Psi\boldsymbol r]=\operatorname{tr}(\Psi\Cov).
\end{equation}
For the exact precision this equals $p$. If the mean is wrong by
$\boldsymbol b$, an additional $\boldsymbol b\T\Psi\boldsymbol b$ appears.
If parameters are fitted or a finite-sample mean is used, the reference
expectation also changes. Do not demand $p$ without checking which residuals
were computed. Examine distributions and parameter-sensitive directions as
well as this single average; different errors can cancel in a trace.

\paragraph{Check the errors that the analysis reports.}
In a locally linear model without priors, let
$H=(D\T\Psi D)^{-1}D\T\Psi$. The fitted parameter displacement is
$H\boldsymbol r$, so its true scatter is $H\Cov H\T$. This need not equal
the quoted covariance $(D\T\Psi D)^{-1}$ when the precision is approximate.
Held-out fits can compare these quantities and test whether, for example,
a nominal 68\% region contains the input parameters at the expected rate.
Priors, nonlinear models and finite validation samples need their own
calibration. Repeating the training/validation split tests sensitivity to
which parents happened to be available; it does not increase the number
of independent parent skies.

\paragraph{What success would buy us.}
A validated correction could retain a full joint vector with more reliable
weights than the analytic model alone, even when the raw sample covariance
cannot be inverted. It could identify covariance directions that limit our
claims about combining two- and three-point statistics. It would not require
compression merely to invert a noisy matrix. Compression remains an option
for a different scientific or computational purpose, with training kept
independent of validation. The immediate deliverables here are the notes,
a two-file viewing example and the PME workflow; measured DES corrections,
resolution convergence and calibrated parameter errors remain future work.

\begin{thebibliography}{9}
\bibitem{pme} O. Friedrich and T. Eifler,
\emph{Precision matrix expansion --- efficient use of numerical simulations
in estimating errors on cosmological parameters},
\href{https://arxiv.org/abs/1703.07786}{arXiv:1703.07786}.
\bibitem{gomes} R. C. H. Gomes et al.,
\emph{Cosmology with second and third-order shear statistics for the Dark
Energy Survey: Methods and simulated analysis},
\href{https://arxiv.org/abs/2503.03964}{arXiv:2503.03964}, version 2.
The public source code is at
\href{https://github.com/rchgomes/3pcf_integrator}{rchgomes/3pcf\_integrator};
the theory work also builds on
\href{https://github.com/git-sunao/fastnc}{git-sunao/fastnc}.
\end{thebibliography}
\end{document}
